\chapter{Rolling Your Own Custom Channel}
\label{chap:custom-channels}

\epigraph{``Give a developer a package, and you feed them for a day. Teach a developer to write a Guix channel, and their entire team never suffers from dependency issues again.''}{--- The Ancient Guild of Guixers}

\noindent
Once you start writing your own Guix package definitions, you will quickly want to share them with your colleagues, your organization, or the broader world.

Rather than emailing \texttt{.scm} files around or maintaining fragile private PPA repositories, GNU Guix allows anyone to turn any Git repository into an authenticated, first-class \textbf{Guix Channel}.

\section{The Anatomy of a Channel Repository}

A custom channel is simply a Git repository structured with standard Scheme module paths. Here is a recommended layout:

\begin{lstlisting}[style=guixcli]
my-guix-channel/
|-- .guix-channel          # Channel metadata file
|-- etc/
|   \-- news.txt           # Announcements displayed on `guix pull`
\-- mychannel/
    |-- packages/
    |   |-- web.scm        # Web applications & services
    |   \-- science.scm    # Scientific computation packages
    \-- services/
        \-- custom.scm     # Custom Shepherd system services
\end{lstlisting}

\section{The \texttt{.guix-channel} Metadata File}

At the root of your repository, create a \texttt{.guix-channel} file:

\begin{lstlisting}[style=guixscheme]
;; .guix-channel
(channel
  (version 0)
  (directory ".")
  ;; Optional: Declare channel dependencies
  (dependencies
    (list (channel
            (name 'guix)
            (url "https://git.savannah.gnu.org/git/guix.git")))))
\end{lstlisting}

\section{Writing Package Modules Inside the Channel}

In \texttt{mychannel/packages/science.scm}, declare your module matching the directory path:

\begin{lstlisting}[style=guixscheme]
(define-module (mychannel packages science)
  #:use-module (guix packages)
  #:use-module (guix download)
  #:use-module (guix build-system gnu)
  #:use-module ((guix licenses) #:prefix license:))

(define-public custom-simulator
  (package
    (name "custom-simulator")
    (version "2.0.0")
    (source (origin
              (method url-fetch)
              (uri "https://internal.example.org/sim-2.0.0.tar.gz")
              (sha256
               (base32
                "0123456789abcdef0123456789abcdef0123456789abcdef0123"))))
    (build-system gnu-build-system)
    (synopsis "Proprietary internal physics simulation engine")
    (description "High performance simulator for organizational workflows.")
    (home-page "https://internal.example.org")
    (license license:gpl3+)))
\end{lstlisting}

\section{Testing Your Channel Locally: \texttt{-L}}

Before committing and pushing your channel to Git, you can test your new packages immediately using the \texttt{-L} (load path) flag:

\begin{lstlisting}[style=guixcli]
# Build the package directly using the local channel directory
guix build -L /path/to/my-guix-channel custom-simulator

# Test running the package in a containerized shell
guix shell -L /path/to/my-guix-channel custom-simulator -- custom-simulator --version
\end{lstlisting}

\section{Channel News: \texttt{etc/news.txt}}

Guix provides a built-in notification system. When users of your channel run \guixcmd{guix pull}, Guix will display formatted news items for new releases or breaking changes:

\begin{lstlisting}[style=guixscheme]
(channel-news
  (version 0)
  (entry (tag "v2.0.0")
         (title (en "Major upgrade: Custom Simulator 2.0 Released!"))
         (body (en "Version 2.0 introduces GPU acceleration and fixes memory leaks."))))
\end{lstlisting}

\begin{protipbox}[Subscribing Teammates to Your Channel]
To enable your team to use your new channel, they simply add this snippet to their \texttt{\textasciitilde/.config/guix/channels.scm}:
\begin{lstlisting}[style=guixscheme]
(cons (channel
        (name 'mychannel)
        (url "https://git.example.org/my-guix-channel.git")
        (branch "main"))
      %default-channels)
\end{lstlisting}
After running \guixcmd{guix pull}, your entire company's internal software library is available directly via \guixcmd{guix install} and \guixcmd{guix shell}!
\end{protipbox}
